\centerline{\bf \Huge Банки. База(А+Д) II}
\centerline{\bf \Huge }

\begin{flushright}
        \scriptsize{\textbf{qq}}
\end{flushright}


\problem{\textbf{1*}} В июле некоторого года планируется взять кредит в банке на некоторую сумму. Условия его возврата таковы:

— каждый январь долг возрастает на 25\% по сравнению с долгом на конец предыдущего
года;

— с февраля по июнь необходимо выплатить часть долга одним платежом.
Известно, что сумма всех выплат составила 375 000 рублей. 

Сколько рублей было взято в банке, если известно, что кредит был полностью погашен четырьмя равными платежами?

\problem{\textbf{2*}} 15 января планируется взять кредит в банке на 25 месяцев. Условия его возврата таковы:

— 1-ого числа каждого месяца долг возрастает на $r$\% по сравнению с концом предыдущего
месяца;

— со 2-го по 14-е число каждого месяца необходимо выплатить часть долга;

— 15-го числа каждого месяца долг должен быть на одну и ту же величину меньше долга
на 15-е число предыдущего месяца.

Известно, что общая сумма денег, которую нужно выплатить банку за весь срок кредитования, на 13\% больше, чем сумма, взятая в кредит. Найдите $r$.

\problem{\textbf{3}} В феврале 2015 года Аркадий Петрович взял кредит в банке под 13\% годовых. Выплатить кредит он должен восемью платежами, вносимыми раз в год на счет после начисления процентов на оставшуюся сумму долга. Долг при этом должен уменьшаться каждый год равномерно.
Сколько рублей составит переплата по кредиту, если наибольший платеж на 91000 рублей
больше наименьшего платежа?

.

.

\problem{\textbf{4}} Банк «Т» предлагает кредит на 3 года на покупку машины стоимостью 546 000
рублей на следующих условиях:

— раз в год банк начисляет на остаток долга 20\%;

— после начисления процентов клиент обязан внести некоторую сумму в счет погашения
части долга;

— выплачивать кредит необходимо равными ежегодными платежами.

Сколько рублей составит переплата по такому кредиту?


\problem{\textbf{5}} Каков процент годовых в банке по кредиту, который выдается на 7 лет, если сумма, выплаченная банку за все годы кредитования, составляет 144\% от суммы кредита, а погашение
кредита происходит раз в год после начисления процентов платежами, уменьшающими долг
равномерно?



\problem{\textbf{6}} Банк выдает кредит сроком на 4 года под 25\% годовых. Вычислите, на сколько процентов переплата по такому кредиту превышает платеж, если гасить кредит нужно равными ежегодными выплатами.

\problem{\textbf{7}} В июле 2025 года планируется взять кредит в банке на некоторую сумму на 10 лет.
Условия его возврата таковы:

— каждый январь долг будет возрастать на 10\% по сравнению с концом предыдущего года;

— с февраля по июнь каждого года необходимо выплатить одним платежом часть долга;

— в июле 2026, 2027, 2028, 2029 и 2030 годов долг должен быть на какую-то одну и ту же
величину меньше долга на июль предыдущего года;

— в июле 2030 года долг должен составлять 800 тыс. рублей;

— в июле 2031, 2032, 2033, 2034 и 2035 годов долг должен быть на другую одну и ту же
величину меньше долга на июль предыдущего года;

— к июлю 2035 года кредит должен быть выплачен полностью.

Известно, что сумма всех платежей после полного погашения кредита будет равна 2090 тыс.
рублей. Какую сумму планируется взять в кредит?

